\section{F. The monodromy group and its period domain}

We use the clockwise root deck generators and the lattice basis specified
on the first page.  In this section a group acts properly discontinuously when, for every
compact subset $K$ of its space, only finitely many group elements $g$ satisfy
$gK\cap K\ne\varnothing$.  This definition permits finite stabilizers.  We
write $e=\widehat\delta$, $E=e\otimes\gamma$, and
\[
 J_2=\begin{pmatrix}0&1\\-1&0\end{pmatrix},\qquad
 \Omega=\begin{pmatrix}
 0&0&0&6\\0&0&1&0\\0&-1&0&0\\-6&0&0&0
 \end{pmatrix}.
\]
Thus the alternating form is
$\psi=u\wedge w+6\gamma\wedge\delta$, an integral alternating form of
type $(1,6)$ on $\Lambda$.  We use $\Omega$ for its matrix and retain $J$ for
a real complex structure.  The domain $\mathcal D$ consists of the real
endomorphisms $J$ of $\Lambda_{\mathbb R}$ satisfying
$J^2=-I$, $J^{\mathsf t}\Omega J=\Omega$, and
$\operatorname{sign}(\Omega J)=(2,2)$.  It is the stated
$\operatorname{Sp}_4(\mathbb R)/U(1,1)$, with its complex structure as an
open subset of the Lagrangian Grassmannian.  An integral symplectic matrix
acts by $J\mapsto gJg^{-1}$.

\subsection{F.1. An exact arithmetic description}

\begin{theorem}[F.1.1; Proved]\label{st:F.1.1}
Let $P_{e,\mathbb Z}$ be the integral symplectic group fixing $e$ as a
vector, and let $P_{\mathbb Z}$ be the integral symplectic group preserving
$\mathbb Ze\subset\ker\gamma\subset\Lambda$.  Every element of
$P_{e,\mathbb Z}$ has a unique expression
\begin{equation}\label{eq:F-g}
 g(B,x,t)=
 \begin{pmatrix}
 1&0&0\\6x&B&0\\t&-x^{\mathsf t}J_2B&1
 \end{pmatrix},\qquad
 B\in\operatorname{SL}_2(\mathbb Z),\quad x\in\mathbb Z^2,\quad t\in\mathbb Z,
\end{equation}
where the block sizes are $1,2,1$.  Put
\[
 S=\begin{pmatrix}0&-1\\1&0\end{pmatrix},\quad
 T=\begin{pmatrix}1&1\\0&1\end{pmatrix},\quad
 a=\begin{pmatrix}0&1\\-1&-1\end{pmatrix}.
\]
There is a homomorphism
$\alpha:\operatorname{SL}_2(\mathbb Z)\to\mathbb Z/12$ determined by
$\alpha(S)=9$ and $\alpha(T)=1$; in particular $\alpha(a)=4$.
With $Q(x)=x_1+x_2+x_1x_2$, the actual monodromy group is
\begin{equation}\label{eq:F-Gamma}
 \Gamma=\left\{g(B,x,t):t+6Q(x)-\alpha(B)=0\ \text{in }\mathbb Z/12\right\}.
\end{equation}
Consequently $[P_{e,\mathbb Z}:\Gamma]=12$ and
$P_{\mathbb Z}=\{\pm I\}\times P_{e,\mathbb Z}$, so
$[P_{\mathbb Z}:\Gamma]=24$.  Its image in the Levi quotient is exactly
$\operatorname{SL}_2(\mathbb Z)$, with scalar on $e$ equal to $1$.
\end{theorem}

\begin{proof}
The matrices in the question, in the fixed basis, are
\[
 A_1=\begin{pmatrix}1&0&0&0\\6&0&1&0\\-6&-1&-1&0\\-2&1&0&1\end{pmatrix},
 \qquad
 A_2=\begin{pmatrix}1&0&0&0\\0&0&-1&0\\-6&1&0&0\\3&0&1&1\end{pmatrix}.
\]
Multiplication gives $A_1^3=A_2^4=I$ and
$A_j^{\mathsf t}\Omega A_j=\Omega$.  In particular
$A_1=g(a,(1,-1)^{\mathsf t},-2)$ and
$A_2=g(S,(0,-1)^{\mathsf t},3)$.  Their images $a,S$ generate
$\operatorname{SL}_2(\mathbb Z)$ because $Sa=T$ and $S,T$ perform the
Euclidean algorithm on primitive integral columns.

A matrix fixing $e$ and preserving $\Omega$ preserves
$e^{\perp}=\ker\gamma$.  Its first row is $(1,0,0,0)$, its last column is
$e$, its middle block $B$ is symplectic on
$e^{\perp}/\mathbb Re$, and its remaining middle column, say $v$, satisfies
$6r=-v^{\mathsf t}J_2B$.  Integrality of $r$ and unimodularity of $B$
give $v=6x$ with $x\in\mathbb Z^2$.  There is no further condition on
the lower left entry $t$.  This proves \eqref{eq:F-g}.  Direct multiplication
gives
\begin{equation}\label{eq:F-law}
 g(B,x,t)g(C,y,s)
 =g(BC,x+By,t+s-6x^{\mathsf t}J_2By).
\end{equation}

For completeness, the presentation of the Levi group needed here follows
from an elementary action on the real projective line.  The images of $a$
and $S$ in $\operatorname{PSL}_2(\mathbb Z)$ have orders $3$ and $2$.
On positive real numbers, $a(z)=-1/(z+1)$ has image $(-1,0)$ and
$a^2(z)=-(z+1)/z$ has image $(-\infty,-1)$; $S$ carries negative real
numbers to positive real numbers.  A cyclically reduced alternating word
starting with a nonidentity power of $a$ and ending with $S$ therefore
carries $(-\infty,0)$ into one of these two proper subintervals.  It cannot
be the identity.  Every reduced word is conjugate to such a word or to a
nonidentity element of one factor.  Together with the Euclidean algorithm,
this proves the presentation
$\operatorname{PSL}_2(\mathbb Z)=\langle\bar a,\bar S\mid
\bar a^3=\bar S^2=1\rangle$.  It follows that
\begin{equation}\label{eq:F-SLpresentation}
 \operatorname{SL}_2(\mathbb Z)=
 \langle a,S\mid a^3=S^4=1,\ [a,S^2]=1\rangle.
\end{equation}
Indeed the group on the right modulo its central element $S^2$ is the
just-proved free product, and $S^2$ maps to the nonidentity matrix $-I$.
Its kernel in the displayed matrix group is consequently trivial.
The values $4,9$ respect all these relations, proving the assertion about
$\alpha$; $T=Sa$ has value $1$.

The quadratic function $Q$ modulo $2$ is zero at $(0,0)$ and one at all
three other points of $\mathbb F_2^2$.  Hence $Q(Bx)=Q(x)\pmod2$ for every
$B\in\operatorname{SL}_2(\mathbb Z)$.  Also
\[
 Q(x+y)-Q(x)-Q(y)\equiv x^{\mathsf t}J_2y\pmod2.
\]
Equation \eqref{eq:F-law} now shows that
$h(g(B,x,t))=t+6Q(x)\pmod{12}$ is a homomorphism.
The homomorphism $\chi=h-\alpha\circ B$ is surjective, since
$\chi(g(I,0,1))=1$.  Its values on $A_1,A_2$ are zero.

We verify equality with its kernel; containment alone would not determine
the group.  Let $N_{\mathbb Z}$ be the kernel of the map to the middle
block.  The kernel of
$\langle g_1,g_2\mid g_1^3=g_2^4=1\rangle\to
\operatorname{SL}_2(\mathbb Z)$ is normally generated by
$[g_1,g_2^2]$, by \eqref{eq:F-SLpresentation}.  Its image in $\Gamma$ is
therefore normally generated by
\[
 K=[A_1,A_2^2]=g(I,(0,-1)^{\mathsf t},6).
\]
Put
\[
 U=A_2KA_2^{-1}=g(I,(1,0)^{\mathsf t},-6),\quad
 V=K^{-1}=g(I,(0,1)^{\mathsf t},-6),\quad Z=I+12E.
\]
Equation \eqref{eq:F-law} gives $[U,V]=Z^{-1}$.  Thus $U,V,Z$ generate
\begin{equation}\label{eq:F-H}
 H=\{g(I,x,t):t\in6\mathbb Z,\quad t/6\equiv Q(x)\pmod2\}.
\end{equation}
More explicitly,
$U^rV^sZ^k=g(I,(r,s)^{\mathsf t},-6(r+s+rs)+12k)$, so these expressions
give every member of the set.  Conjugation by $g(B,x,t)$ sends
$g(I,y,s)$ to $g(I,By,s-12x^{\mathsf t}J_2By)$; the invariance of $Q$
modulo $2$ proves that $H$ is normal in $P_{e,\mathbb Z}$.  The normal
closure of $K$ is contained in $H$ and contains $U,V,Z$, so it is exactly
$H$.  We have proved $\Gamma\cap N_{\mathbb Z}=H$.  Both $\Gamma$ and
$\ker\chi$ surject onto the same Levi group and have intersection $H$
with $N_{\mathbb Z}$.  Given $g\in\ker\chi$, choose $\gamma\in\Gamma$
with the same middle block.  Then $g\gamma^{-1}\in H$, and $g\in\Gamma$.
This proves \eqref{eq:F-Gamma} and the index $12$.

Finally an integral automorphism preserving the primitive line
$\mathbb Ze$ sends $e$ to $e$ or $-e$.  Multiplication by the central
matrix $-I$ converts the latter case to the former.  The intersection of
$\{\pm I\}$ with $P_{e,\mathbb Z}$ is trivial.  This proves the direct
product and the index $24$.
\end{proof}

\begin{proposition}[F.1.2; Proved]\label{st:F.1.2}
The group in \eqref{eq:F-Gamma} has center
$\langle Z\rangle\cong\mathbb Z$ and is the following explicit central
extension of $\mathbb Z^2\rtimes\operatorname{SL}_2(\mathbb Z)$ by
$\mathbb Z$.  Choose the representative $a(B)\in\{0,\ldots,11\}$ of
$\alpha(B)$.  Its underlying set consists of triples $(B,x,k)$, with
product
\[
 (B,x,k)(C,y,l)=(BC,x+By,k+l+\kappa((B,x),(C,y))),
\]
where
\begin{align*}
12\kappa((B,x),(C,y))={}&a(B)+a(C)-a(BC)-6Q(x)-6Q(y)\\
 &-6x^{\mathsf t}J_2By+6Q(x+By).
\end{align*}
Its real and rational Zariski closures are respectively the full real
and rational symplectic stabilizers of the vector $e$.  They have
dimension $6$, with three-dimensional Heisenberg unipotent radical and
Levi quotient $\operatorname{SL}_2$.
\end{proposition}

\begin{proof}
The numerator defining $\kappa$ is divisible by $12$, by the two
congruences proved above.  The bijection
\[
 (B,x,k)\longmapsto g(B,x,a(B)-6Q(x)+12k)
\]
intertwines the displayed multiplication with matrix multiplication by
\eqref{eq:F-law}.  This proves associativity as well as the claimed abstract
description, without leaving an unspecified extension class.  The kernel
of the homomorphism to $\mathbb Z^2\rtimes\operatorname{SL}_2(\mathbb Z)$
is $\langle Z\rangle$.  That affine group has trivial center: commuting
with all translations forces its linear part to be $I$, and commuting
with $S,T$ forces its translation part to vanish.  Thus the center of
$\Gamma$ is exactly the stated kernel.

The logarithms $U-I,V-I,Z-I$ are linearly independent elements of the
three-dimensional Heisenberg Lie algebra.  The square of each is zero.
A polynomial vanishing on all integral powers $I+n(U-I)$ vanishes on
the whole one-parameter line; the same holds for $V$ and $Z$.  Hence the
Zariski closure contains the full unipotent radical.  Its Levi image
contains $S,T$; the powers of $T$ and of $STS^{-1}$ similarly supply the
two opposite unipotent subgroups of $\operatorname{SL}_2$.  Their Lie
algebras and their bracket span $\mathfrak{sl}_2$.  They generate the
connected algebraic group $\operatorname{SL}_2$.  A subgroup containing
the entire unipotent kernel and surjecting onto this Levi quotient is
the entire vector stabilizer.  Its dimension is $3+3=6$.  This reasoning
uses equations over $\mathbb Q$, so proves the rational assertion as
well as the real assertion.
\end{proof}

\subsection{F.2. Properness and the maximal open set}

Define the real-analytic function
\[
 q_e:\mathcal D\longrightarrow\mathbb R,\qquad q_e(J)=\psi(e,Je),
 \qquad\mathcal O=\{J:q_e(J)\ne0\}.
\]
The function $q_e$ is invariant under $P_{\mathbb Z}$: replacing $e$ by
$-e$ changes both of its arguments.  Write $\mathcal O_+$ and
$\mathcal O_-$ for its two sign regions.

\begin{proposition}[F.2.1; Refuted]\label{st:F.2.1}
The assertion that $\Gamma$ acts properly discontinuously on all of
$\mathcal D$ is false.  The following point of $\mathcal D$ has an
explicit infinite stabilizer:
\[
 J_0=\begin{pmatrix}
 0&0&-1&0\\0&0&0&-6\\1&0&0&0\\0&1/6&0&0
 \end{pmatrix},\qquad
 R=\begin{pmatrix}
 1&0&0&0\\0&1&72&0\\0&0&1&0\\-12&0&0&1
 \end{pmatrix}\in\Gamma.
\]
Every power of $R$ fixes $J_0$, and all these powers are distinct.
\end{proposition}

\begin{proof}
Direct multiplication gives $J_0^2=-I$ and
$J_0^{\mathsf t}\Omega J_0=\Omega$.  Moreover
\[
 \Omega J_0=\begin{pmatrix}
 0&1&0&0\\1&0&0&0\\0&0&0&6\\0&0&6&0
 \end{pmatrix}
\]
has eigenvalues $1,-1,6,-6$, so $J_0\in\mathcal D$ and $q_e(J_0)=0$.
The matrix $R$ equals $g(T^{72},0,-12)$ and belongs to \eqref{eq:F-Gamma},
because $\alpha(T^{72})=0\pmod{12}$.  Its commutator with $J_0$ is zero.
Writing $R=I+N$ gives $N\ne0$ and $N^2=0$, so $R^m=I+mN$ and the
powers are distinct.  The compact singleton $\{J_0\}$ therefore violates
the stated definition of proper discontinuity.  These identities are
checked both with SymPy and with a separate rational matrix
implementation in \texttt{code/f\_exact.py}.
\end{proof}

\begin{theorem}[F.2.2; Proved]\label{st:F.2.2}
The set $\mathcal O$ is the largest $\Gamma$-invariant open subset of
$\mathcal D$ on which the action is properly discontinuous.  It is also
the largest such open subset for $P_{\mathbb Z}$.  Each sign region is
a complex manifold in the analytic topology, and its quotient by either
discrete group is a Hausdorff complex analytic orbifold.
\end{theorem}

\begin{proof}
Let $P_{\mathbb R}$ be the real symplectic group preserving the real line
$\mathbb Re$.  We first prove properness on each sign region.  For a fixed
$J$, the form
$h_J(v,w)=\psi(v,Jw)+i\psi(v,w)$ is a nondegenerate Hermitian form of
signature $(1,1)$ on the complex vector space $(\Lambda_{\mathbb R},J)$,
with the convention conjugate linear in the first entry.  When
$q_e(J)\ne0$, the complex line $\mathbb Ce$ is nondegenerate and its
Hermitian orthogonal complement is a complex line of the opposite sign.
Normalize $e$ by a positive real scalar and choose a unit vector in this
complement.  Comparing these two Hermitian bases for two complex
structures with the same sign gives a real symplectic map preserving
$\mathbb Re$ and conjugating one complex structure to the other.  Thus
$P_{\mathbb R}$ is transitive on each sign region.

An element of the stabilizer of $J$ inside $P_{\mathbb R}$ is complex
linear and sends $e$ to $r e$ for a real nonzero $r$.  Preservation of
$q_e(J)$ gives $r^2=1$.  On the complementary complex line it acts by an
element of $U(1)$.  The stabilizer is consequently the compact group
$\{\pm1\}\times U(1)$.  Each sign region is therefore a homogeneous
space of $P_{\mathbb R}$ with compact stabilizer $L$.  Local sections of
$P_{\mathbb R}\to P_{\mathbb R}/L$, and a finite cover of a compact
subset $K$, give a compact subset $C$ of $P_{\mathbb R}$ mapping onto
$K$.  Every element carrying a point of $K$ back into $K$ belongs to
$CLC^{-1}$, which is compact.  This proves properness.  Restricting a
proper action to a closed discrete subgroup gives
the required finite-intersection condition.  Both $\Gamma$ and
$P_{\mathbb Z}$ are closed discrete matrix groups.  The two sign regions
are disjoint invariant open-and-closed subsets of $\mathcal O$, so the
same finite-intersection condition holds on their union.  Proper actions
by discrete groups have finite stabilizers and Hausdorff quotients; their
holomorphic local finite-group quotient charts supply the asserted
analytic orbifolds.

It remains to prove maximality, which requires more than one bad point.
The real vector stabilizer $P_{e,\mathbb R}$ acts transitively on
$\mathcal D_0=\{J:q_e(J)=0\}$.  Here is a direct comparison map.  In the
Hermitian space just described, $e$ is a nonzero null vector.  Choose
$v$ with $h_J(e,v)=1$, and replace $v$ by
$v-\tfrac12h_J(v,v)e$.  This preserves $h_J(e,v)=1$ and makes $v$ null.
The ordered complex basis $(e,v)$ has Gram matrix
$\bigl(\begin{smallmatrix}0&1\\1&0\end{smallmatrix}\bigr)$.
The complex-linear map between two such bases fixes $e$ and preserves
$h$, hence its imaginary part $\psi$.  As a real map it belongs to
$P_{e,\mathbb R}$ and conjugates the two complex structures.

The subgroup $P_{e,\mathbb Q}$ is dense in $P_{e,\mathbb R}$.
For the unipotent coordinates in \eqref{eq:F-g} this is density of
$\mathbb Q^3$ in $\mathbb R^3$.  For $\operatorname{SL}_2$ one may
approximate $a,b,c$ by rational numbers in the chart $a\ne0$ and set
$d=(1+bc)/a$; translates of that chart cover the group.  Hence
\[
 \{rJ_0r^{-1}:r\in P_{e,\mathbb Q}\}
\]
is dense in $\mathcal D_0$.
Put $R_0=g(T^6,0,-1)$, so $R_0^{12}=R$ and
$(R_0-I)^2=0$.  At $rJ_0r^{-1}$ the nonzero rational nilpotent
$L=r(R_0-I)r^{-1}$ commutes with the complex structure and satisfies
$L^2=0$.  Choose a positive integer $d$ clearing its denominators.
The matrix $I+dL$ and its inverse $I-dL$ are integral, symplectic, and
fix $e$; this follows either by conjugating the identity
$(I+s(R_0-I))^{\mathsf t}\Omega(I+s(R_0-I))=\Omega$ or by expanding it.
The twelfth power of $I+dL$ lies in $\Gamma$, because
$P_{e,\mathbb Z}/\Gamma\cong\mathbb Z/12$.  It has infinite order and
fixes $rJ_0r^{-1}$.  Thus points with infinite $\Gamma$-stabilizer are
dense in $\mathcal D_0$.

Every open neighborhood in $\mathcal D$ of any point of $\mathcal D_0$
therefore contains a point with infinite stabilizer.  An open set on
which the action is proper cannot meet $\mathcal D_0$.  This proves the
largest-open-set assertion for $\Gamma$.  The same points have infinite
$P_{\mathbb Z}$-stabilizer, so it proves that assertion for the full
parabolic as well.
\end{proof}

\begin{proposition}[F.2.3; Proved]\label{st:F.2.3}
The positive sign region admits the global period chart
\[
 \mathcal O_+=\left\{
 (\tau,\mu,\beta)\in\mathbb C^3:
 \operatorname{Im}\tau>0,\quad
 \operatorname{Im}\beta-\frac{6(\operatorname{Im}\mu)^2}
 {\operatorname{Im}\tau}<0\right\},
\]
represented by the matrix $\Pi$ in the question.  In this chart,
\begin{equation}\label{eq:F-qD}
 q_e(J)=-\frac6D.
\end{equation}
Every admissible period $\Pi_c(z)$ belongs to $\mathcal O_+$, for every
admissible $c$.  Thus this single maximal proper-action open set contains
all the period images in question.
\end{proposition}

\begin{proof}
Let $J\in\mathcal O_+$.  The two vectors $\widehat w,e$ are complex-linearly
independent in $(\Lambda_{\mathbb R},J)$.  Otherwise
$\widehat w=r e+sJe$ for real $r,s$.  Pairing with $e$ gives
$0=sq_e(J)$; hence $s=0$, contradicting the real linear independence
of $\widehat w,e$.  There is consequently a unique complex-linear
isomorphism $(\Lambda_{\mathbb R},J)\to\mathbb C^2$ sending these two
vectors to $(1,0)$ and $(0,1)$.  The corresponding period matrix has
right block $I$.  The equation
$\Pi\Omega^{-1}\Pi^{\mathsf t}=0$ says that its upper left entry is
six times its lower middle entry, giving exactly the displayed matrix.
This construction is the ordinary holomorphic Grassmannian chart,
so its coordinates and inverse are holomorphic.

In the real symplectic basis
$(\widehat u,\widehat\gamma/6,\widehat w,e)$ the normalized symmetric
period matrix is
$\bigl(\begin{smallmatrix}\tau&\mu\\\mu&\beta/6\end{smallmatrix}\bigr)$.
Its imaginary part has signature $(1,1)$, equivalently
$6(\operatorname{Im}\mu)^2-
\operatorname{Im}\tau\operatorname{Im}\beta>0$.
Solving $\Pi(Je)=i\Pi e=(0,i)$ in real coefficients gives
\[
 \gamma(Je)=
 \frac{\operatorname{Im}\tau}
 {\operatorname{Im}\tau\operatorname{Im}\beta-
 6(\operatorname{Im}\mu)^2}.
\]
Since $\psi(e,v)=-6\gamma(v)$, positivity of $q_e$ now gives
$\operatorname{Im}\tau>0$ and $D<0$, and proves \eqref{eq:F-qD}.
Conversely these inequalities make the imaginary period matrix
invertible with the stated signature and give $q_e>0$.  The admissibility
inequality in the question is $D_c<0$ everywhere, which proves the last
assertion by this exact comparison map.
\end{proof}

\subsection{F.3. The curve in the proper quotient}

\begin{theorem}[F.3.1; Proved]\label{st:F.3.1}
For every admissible $c$, the marked period map on the covering of
$B^\circ$ is a holomorphic immersion into $\mathcal O_+$.  The induced
map
\[
 p_c:B^\circ\longrightarrow\Gamma\backslash\mathcal O_+
\]
is injective and a holomorphic immersion.  It extends over $t=0,1$ to a
proper holomorphic embedding
\[
 \bar p_c:\mathbb C=B\setminus\{\infty\}
 \lhook\joinrel\longrightarrow\Gamma\backslash\mathcal O_+.
\]
The closure of $p_c(B^\circ)$ in this quotient is exactly
$\bar p_c(\mathbb C)$.
\end{theorem}

\begin{proof}
For $B=\bigl(\begin{smallmatrix}a&b\\c&d\end{smallmatrix}\bigr)$, changing
the marking by $g(B,x,t)$ replaces $\Pi$ by $\Pi g(B,x,t)^{-1}$ and then
renormalizes its right block.  On the quotient by the complex line of
$e$, the two periods are the row $(\tau,1)$.  Right multiplication by
$B^{-1}$ sends it to $(d\tau-c,a-b\tau)$.  Hence
\begin{equation}\label{eq:F-tauaction}
 \tau\longmapsto\frac{d\tau-c}{a-b\tau}.
\end{equation}
The denominator never vanishes in the upper half plane.  The function
$j(\tau)$ is therefore a holomorphic invariant of this action and
descends to a holomorphic function $j_\Gamma$ on the proper quotient.
Here $j$ is the elliptic invariant, normalized by
$j(E)=1728\,4A^3/(4A^3+27B^2)$ for $y^2=x^3+Ax+B$; changing a lattice
basis preserves the isomorphism of elliptic curves and hence this
function.

We derive the required value in the present family.  The elliptic curve
whose period is $\tau$ has the equation stipulated in the question,
\[
 y^2=x^3-3t^3(t-1)x+2t^4(t-1)^2.
\]
Its discriminant is $1728t^8(t-1)^3$ and direct cancellation gives
\begin{equation}\label{eq:F-j}
 j(\tau)=1728t\qquad(t\in B^\circ).
\end{equation}
Equivalently this is the defining lift of the elliptic modular map used
for $\tau$ in the period construction.  The computation is reproduced
exactly in \texttt{f\_exact.py}.  In a local base coordinate $t$, the
derivative of the composite $j\circ\tau$ equals $1728$.  Thus the
derivative of the marked period map is nonzero.  The same identity after
passage to the quotient shows both injectivity and injectivity of the
tangent map of $p_c$: $j_\Gamma/1728$ is a holomorphic left inverse.

At the two elliptic points, the period functions are holomorphic in the
cyclic covering coordinates $s_1,s_2$, with $t=s_1^3$ and
$t-1=s_2^4$ after choosing local coordinates as in the question.  Their
values are nondegenerate because $D_c<0$ there.  The maps from these
discs to $\mathcal O_+$ are equivariant for $A_1,A_2$.  Composing with
the quotient makes them invariant, so they descend to holomorphic maps
from the base discs.  This is the finite-quotient comparison map: a
holomorphic invariant germ on a disc is a convergent series in $s_j^{m_j}$,
which proves descent on local function rings.  The extensions agree
with $p_c$ off their centers and therefore with each other.  Identity
\eqref{eq:F-j} extends by holomorphy, so the resulting $\bar p_c$ still
has holomorphic left inverse $j_\Gamma/1728$.

The quotient is Hausdorff by F.2.2.  For a compact subset $K$ of it,
$\bar p_c^{-1}(K)$ is closed in $\mathbb C$ and lies in the compact set
$j_\Gamma(K)/1728$.  Thus $\bar p_c$ is proper.  A holomorphic section of
a holomorphic map is locally a closed immersion: in local embeddings
into polydiscs, its coordinate along the base is $t$ and the remaining
coordinates are its holomorphic component functions, whose graph
defines the image.  Applying this to the left inverse shows that
$\bar p_c$ is a holomorphic embedding; properness makes its image
closed.  Since $\mathbb C\setminus\{0,1\}$ is dense in $\mathbb C$, its
image has the asserted closure.
\end{proof}

\begin{proposition}[F.3.2; Proved]\label{st:F.3.2}
At the two finite punctures the period map has the holomorphic
finite-cover limits
\[
 \Pi_1=
 \begin{pmatrix}2(2-\rho)&\rho&1&0\\
 \beta_c(z_1)&(2-\rho)/3&0&1\end{pmatrix},\qquad
 \Pi_2=
 \begin{pmatrix}3(1-i)&i&1&0\\
 \beta_c(z_2)&(1-i)/2&0&1\end{pmatrix},
 \quad \rho=e^{\pi i/3}.
\]
Their monodromies are the given order-three and order-four matrices.
At the cusp the curve leaves every compact subset of
$\Gamma\backslash\mathcal O_+$.  On the prescribed cusp branch its
Hodge plane in the compact dual tends to
$\mathbb C\gamma+\mathbb Cu$.  The latter plane has complex dimension
two in common with its conjugate and therefore does not represent a
point of $\mathcal D$.
\end{proposition}

\begin{proof}
The finite values of $\tau,\mu$ follow by evaluating their transformation
laws at the fixed points: $\tau(z_1)=\rho$,
$\mu(z_1)=(2-\rho)/3$, $\tau(z_2)=i$,
$\mu(z_2)=(1-i)/2$.  Substitution in $\Pi$ gives the two matrices.
The finite covering comparison and its descent were proved in F.3.1.
The unipotent cusp transformation is
$\tau\mapsto\tau-1$, $\mu\mapsto\mu$, $\beta\mapsto\beta+1$.
On the cusp branch, $\operatorname{Im}\tau\to+\infty$,
$\mu\to\mu_0$, and $\beta_c+\tau\to b_0+c$, from the holomorphic
cusp expansions.  Consequently
\[
 D_c=-\operatorname{Im}\tau+O(1),\qquad
 q_e(\Pi_c)\longrightarrow0^+,
 \qquad j_\Gamma(p_c(t))=1728t\longrightarrow\infty.
\]
The last equality alone excludes a convergent subnet in any compact
subset of the proper quotient; the first two identities give the
position relative to the maximal proper-action open set.
The rows of $\Pi_c$, regarded as covectors, are
$6\mu\gamma+\tau u+w$ and $\beta_c\gamma+\mu u+\delta$.
Dividing each by $\tau$ gives limits $u$ and $-\gamma$.
Their exterior product has a nonzero limit, so their span converges
in the Grassmannian to the stated plane.  It is spanned by real
covectors, hence equals its conjugate.  A complex structure has
$H^{1,0}\cap H^{0,1}=0$; here the intersection is the entire
two-dimensional plane.  This explicit comparison proves the assertion
about the compact-dual limit.
\end{proof}

\begin{proposition}[F.3.3; Proved]\label{st:F.3.3}
The determinant used in the workbench is generically nonzero, with
the stated zero at $z_1$:
\[
 \det\dot Z=6\dot\mu^2-\dot\tau\dot\beta,\qquad
 \det\dot Z(z_1)=0.
\]
The immersion established in F.3.1 holds at every point of $B^\circ$
without requiring this determinant to be everywhere nonzero.
\end{proposition}

\begin{proof}
In the cusp coordinate $s$, the expansion
$Z=s\bigl(\begin{smallmatrix}0&1\\-1&0\end{smallmatrix}\bigr)+C(q)$,
$q=e^{2\pi is}$, gives
$\det(dZ/ds)=1+O(q)$, so the determinant is not the zero function.
At $z_1$ the fixed-point transformation law of $\mu$ gives
$\dot\mu=-\dot\tau/3$.  One may obtain this by differentiating
$\mu\circ g_1=(1-\mu)/\tau$, using
$g_1'(z_1)=-\rho$ and $\rho^2-\rho+1=0$.
The law
$\beta\circ g_1=\beta+2-6(1-\mu)^2/\tau$ then gives
$\dot\beta=2\dot\tau/3$.  Substitution yields
$6\dot\tau^2/9-2\dot\tau^2/3=0$.
The derivative of a map from a curve to the three-dimensional period
chart is a vector of three coordinate derivatives.  Identity
\eqref{eq:F-j} proves that this vector is nonzero over $B^\circ$;
the determinant tests the rank of the associated two-by-two
Kodaira--Spencer matrix.  The maps compared here are explicitly the
coordinate derivative and its determinant, and the displayed computation
states exactly what each supplies.
\end{proof}
